\section{Model computation and sparse-interface background}
\label{app:model-background}

This section introduces the model-side objects needed to state our problem. We first describe how a vision--language model (VLM) produces an answer and what its dense internal state represents. We then introduce sparse transcoders as separate tools for viewing part of the dense internal state. With the base computation and the fitted measurement interface separated, we can state which internal objects may preserve the answer effect of localized visual evidence.

\stitle{VLM computation.}
A VLM receives an image and a question and generates an answer one token at a time. Answer generation is autoregressive: after producing or receiving an answer prefix, the model uses that prefix together with the image and question to predict the next token. We write $x=(I,q)$ for the input image $I$ and question $q$, and $a=(a_1,\ldots,a_T)$ for a fixed target answer containing $T$ tokens. The index $t\in\{1,\ldots,T\}$ identifies one answer token, $a_{<t}$ is the prefix preceding $a_t$, and $\mathcal{V}$ is the set of tokens in the model vocabulary. The frozen parameters of the VLM are denoted by $\theta$.

Before the model predicts $a_t$, it represents the image, question, and answer prefix as one joint sequence of visual and text positions. This sequence passes through $L$ residual-stream layers indexed by $\ell\in\{0,\ldots,L-1\}$. At every position, the residual stream is the running internal vector that carries information from one layer to the next while the layer's attention and feed-forward calculations update it \citep{elhage2021transformerCircuits}. Attention combines information across positions, whereas the feed-forward calculation transforms the vector at each position. Let $n_{\mathrm{pos}}$ be the number of positions in the joint sequence and $d$ the dimension of each residual-stream vector. We write
\begin{equation}
    \vec{h}_{i}^{\ell}(x,a_{<t};\theta)\in\mathbb{R}^{d}
    \label{eq:dense-state}
\end{equation}
for the residual-stream state at position $i\in\{1,\ldots,n_{\mathrm{pos}}\}$ in layer $\ell$. Because this vector retains all $d$ coordinates at that position, we refer to it as a \emph{dense} state.

After the final layer, the model produces a vector $\vec{z}_{t}(x,a_{<t};\theta)\in\mathbb{R}^{|\mathcal{V}|}$. Each entry is the model's score for one possible next token before those scores are converted to probabilities; these entries are called logits. The entry associated with $a_t$ therefore records the model's direct support for the next target token. Repeating this prediction step generates the complete answer.

\stitle{Sparse representations.}
The dense state records the full internal vector at one layer and position, but its individual coordinates do not automatically correspond to separately understandable parts of the computation. A sparse transcoder provides another view. It is trained after the base VLM and kept frozen during analysis, so it is a measurement tool rather than a component that the VLM itself uses to generate the answer. The tool reads a dense state and describes a change that a layer adds to the dense state using a larger set of learned features, while allowing only a small number of those features to be active at once. This limited feature activity is what makes the representation sparse.

Two forms of transcoder are used in this paper. A per-layer transcoder (PLT) uses features to approximate the update produced by the feed-forward calculation within one layer \citep{dunefsky2024transcoders}. A cross-layer transcoder (CLT) reads a feature in one layer and allows that feature to describe contributions written into later layers \citep{ameisen2025circuittracing}. Applying these interfaces to VLM computation follows recent multimodal circuit-tracing work \citep{yang2026vlmcircuittracing}. The distinction matters because a PLT feature has a within-layer target, whereas a CLT feature can have different target layers.

\begin{samepage}
Let $\phi$ denote the frozen parameters of one model-specific sparse interface. In read layer $\ell$, its encoder converts a dense state into a sparse activation vector,
\begin{equation}
    \vec{f}_{i}^{\ell}(x,a_{<t};\theta,\phi)
    =\mathtt{Encode}^{\ell}\!\left(\vec{h}_{i}^{\ell}(x,a_{<t};\theta);\phi\right).
    \label{eq:sparse-encoding}
\end{equation}
This vector lies in $\mathbb{R}^{n_{\mathrm{feat}}^{\ell}}$, where $n_{\mathrm{feat}}^{\ell}$ is the number of available sparse features in layer $\ell$. The index $j\in\{1,\ldots,n_{\mathrm{feat}}^{\ell}\}$ identifies one feature, and its scalar activation is $f_{i,j}^{\ell}=\bigl[\vec{f}_{i}^{\ell}\bigr]_{j}$. Each active feature is paired with a decoder direction that maps its contribution back into the VLM's dense vector space. The Method defines how features are selected, how PLT and CLT decoder directions are indexed by their target layers, and how selected features are reconstructed. A feature index belongs to one fitted interface and does not identify the same internal component in another model or interface.
\end{samepage}

\stitle{Problem definition.}
A \emph{human-localized answer-bearing region} is an image region that an annotator identifies as containing visual information needed for the target answer. Given an image--question input, a target answer, a frozen VLM, and such a region, our first question is whether changing that region alters the model's score for the fixed answer more than changing comparable non-evidence content. Our second question is which internal object retains the answer-relevant internal change caused by the regional perturbation: the complete dense state, the part reconstructed by a fitted sparse interface, selected feed-forward channels that belong to the base model, or a low-dimensional subspace shared across examples.

These questions are related but not equivalent. Showing that changing the complete dense state changes the target-answer score does not show that a sparse reconstruction, a channel group, or a shared low-dimensional projection reproduces the same score change. Likewise, observing different sparse patterns across models does not by itself show that the underlying models use different complete mechanisms, because each restricted object may retain a different part of the dense computation. The next section defines the intervention protocol used to separate these questions.
